\section{External comparison}
\label{sec:external}

Published work enters this program as evidence rather than as authority.
A paper is unchecked until something checks it, exactly like an internal
document; the reason an external agreement counts for more is not that
it is published but that it was produced without knowledge of this
program, so agreement is not bookkeeping.

The map holds \LitPapers{} indexed papers and \LitEdges{} recorded
relations to named claims here. \LitVerified{} of those edges are
verified --- the paper was obtained, read, and compared against the
claim --- across \LitObtained{} distinct papers. The remainder are
recorded targets for papers not yet obtained, which is itself useful: a
paper heavily cited within this web and still unobtained is
automatically the next acquisition target.

An entry earns its place by naming a specific claim and saying what
relationship it has to it. \texttt{workhouse lit --for G19} answers the
query in the other direction, and validation rejects an entry whose
target does not resolve.

\subsection{A 1989 table, and rationals computed cold}

The sharpest external check available to this program is a
Hamiltonian strong-coupling series for $\SU(3)$ glueball masses
published in 1989~\cite{HAMER_1989}, whose Table~1 is pinned here by
digest. Under the normalization bridge $m_n = 2^{\,n-1} a_n$, its
coefficients are compared against the $\Gamma$-point coefficients this
program derives as exact rationals.

In the $C$-odd ($1^{+-}$) channel --- the sector of the carrier, whose
series is \eqref{eq:thirdorder} --- the agreement runs through third
order:

\begin{center}
\begin{tabular}{@{}l l l l@{}}
\toprule
Order & Published & This program & Gap (bound) \\
\midrule
$n=0$ & $\tfrac{16}{3}/2$ & $\tfrac83 = E_{\mathrm{flat}}(0)$ & exact \\
$n=1$ & $a_1 = +1$ & the $u$ coefficient $+1$ & exact \\
$n=2$ & $2a_2 = 0.0359477124184$ & $\tfrac{11}{306}$
      & $9.94\times10^{-14}$ ($2\times10^{-13}$) \\
$n=3$ & $4a_3 = -0.437135556836$ & $-\tfrac{109151}{249696}$
      & $1.11\times10^{-12}$ ($3\times10^{-12}$) \\
\bottomrule
\end{tabular}
\end{center}

\noindent In the $C$-even ($0^{++}$) channel, $2a_2 = -0.21274509804$
against $-\tfrac{217}{1020}$ with gap $7.84\times10^{-13}$, and
$4a_3 = -0.1039004229144$ against $-\tfrac{54049}{520200}$ with gap
$1.36\times10^{-13}$. The sign structure matches as well: $a_1 = -1$ in
the scalar channel, where the carrier's is $+1$.

Each bound above is the printed half-ulp of the published table scaled
with margin, so these are agreements to the precision at which the 1989
table was printed --- not approximate corroborations. The rationals
$11/306$ and $109151/249696$ are the third-order coefficients of
$E_{\mathrm{flat}}$ in \eqref{eq:thirdorder}, computed here from Haar
integration and the cellular complex with no knowledge of the table.

At fourth order, $8a_4 = -0.7751458630184$ against
$m^{(4)}_{\Gamma} = -0.7751458630189173$, gap $5.17\times10^{-13}$
against a bound of $5.3\times10^{-13}$.

\begin{scopenote}
These are \Ttwo{} checks: float agreement within a stated tolerance, to
the precision of a printed table. They corroborate the coefficients;
they do not prove them, and the exact statements of
\S\ref{sec:second}--\S\ref{sec:grading} do not rest on them. Note also
that the pinned revision of the source compares only $2a_2$ and $4a_3$
in its own text --- the fourth-order agreement is registered here and
not leaned on there.
\end{scopenote}

\subsection{The two contradicting edges}

Two edges in the map are recorded as \emph{contradicts}, both from a
1976 strong-coupling series, and both against this program: one against
the third-order $C$-odd coefficient, one against the $\Gamma$-point
fourth-order value. They originate in a disagreement the 1989 paper
itself reports at orders $x^3$ and $x^4$ with the series it supersedes.

Both are recorded as not yet obtained. Until that primary source is
read, this program has two unresolved external contradictions against
coefficients it otherwise regards as settled, and says so rather than
resolving them by the age of the sources. Obtaining that paper is
listed with the acquisition targets below.

\subsection{How far a novelty claim here can be trusted}

Not as far as the size of the index suggests, and the limitation is
worth stating precisely rather than as a disclaimer.

Only \LitVerified{} of \LitEdges{} recorded edges are verified. The
unverified remainder is not a random sample: it includes several of the
\emph{nearest} neighbours to this program's own results. The main
Nuclear Physics B paper of the 1981 strong-coupling glueball series is
recorded as not obtained --- only its errata were --- and the 1974 and
1979 large-$N$ papers that bear most directly on the rank grading of
\S\ref{sec:grading} are likewise unobtained, with the ledger noting for
one of them that its definitions were used from memory and only their
sign consequences checked.

Consequently: where this document says a result is novel, that means
\emph{no verified entry in this index contains it}, which is a weaker
statement than priority. Two of the strongest results are affected. The
rank-grading cancellation of \S\ref{sec:grading} sits closest to
large-$N$ counting arguments in papers this program has not read; and
the fixed-spacing transfer construction of \S\ref{sec:wilson} sits
alongside the classical positivity and self-adjoint transfer-matrix
results of Osterwalder and Seiler, which the index does register.

Acquiring the unobtained near neighbours is cheap relative to
everything else in \S\ref{sec:obligations}, and it is the one item on
that list that could \emph{subtract} from the program rather than add
to it. That is a reason to do it, not a reason to defer it.

% Generated by paper/master/generate_apparatus.py.  Do not edit.
\begin{longtable}{@{}l l l l l@{}}
\caption{Verified external-evidence edges. Each row is a published result that was obtained, read, and checked against a named claim of this program. Relation is the ledger's own classification; \emph{contradicts} and \emph{confusable} are listed first because they are the rows that cost the program something.}\label{tab:literature}\\
\toprule
Paper & Ref. & Year & Bears on & Relation \\
\midrule
\endfirsthead
\toprule Paper & Ref. & Year & Bears on & Relation \\ \midrule \endhead
\bottomrule
\endfoot
\texttt{\small HSB\_\allowbreak{}2000} & \cite{HSB_2000} & 2000 & \texttt{HAMER\_\allowbreak{}A4\_\allowbreak{}NUM} & confusable with \\
\texttt{\small HAMER\_\allowbreak{}1989} & \cite{HAMER_1989} & 1989 & \texttt{BAND\_\allowbreak{}EVEN\_\allowbreak{}BOTTOM} & corroborates \\
\texttt{\small HAMER\_\allowbreak{}1989} & \cite{HAMER_1989} & 1989 & \texttt{C1} & corroborates \\
\texttt{\small LLL\_\allowbreak{}2006} & \cite{LLL_2006} & 2006 & \texttt{C1} & corroborates \\
\texttt{\small HAMER\_\allowbreak{}1989} & \cite{HAMER_1989} & 1989 & \texttt{D\_\allowbreak{}3} & corroborates \\
\texttt{\small SCHIERHOLZ\_\allowbreak{}1988} & \cite{SCHIERHOLZ_1988} & 1988 & \texttt{G18} & corroborates \\
\texttt{\small HAMER\_\allowbreak{}1989} & \cite{HAMER_1989} & 1989 & \texttt{M3\_\allowbreak{}EVEN\_\allowbreak{}K0} & corroborates \\
\texttt{\small KPS\_\allowbreak{}1981} & \cite{KPS_1981} & 1981 & \texttt{R14} & corroborates \\
\texttt{\small CBB\_\allowbreak{}2026} & \cite{CBB_2026} & 2026 & \texttt{R2} & corroborates \\
\texttt{\small KPS\_\allowbreak{}1981} & \cite{KPS_1981} & 1981 & \texttt{G7} & supplies a value \\
\texttt{\small HAMER\_\allowbreak{}1989} & \cite{HAMER_1989} & 1989 & \texttt{HAMER\_\allowbreak{}A4\_\allowbreak{}NUM} & supplies a value \\
\texttt{\small MUNSTER\_\allowbreak{}1985\_\allowbreak{}TM} & \cite{MUNSTER_1985_TM} & 1985 & \texttt{C2} & supplies a comparison \\
\texttt{\small CAO\_\allowbreak{}2023} & \cite{CAO_2023} & 2023 & \texttt{C7} & supplies a comparison \\
\texttt{\small LEMOINE\_\allowbreak{}2026} & \cite{LEMOINE_2026} & 2026 & \texttt{C7} & supplies a comparison \\
\texttt{\small HAZRA\_\allowbreak{}2026} & \cite{HAZRA_2026} & 2026 & \texttt{G14} & supplies a comparison \\
\texttt{\small AT\_\allowbreak{}2020\_\allowbreak{}SU3} & \cite{AT_2020_SU3} & 2020 & \texttt{G18} & supplies a comparison \\
\texttt{\small BESIII\_\allowbreak{}2026} & \cite{BESIII_2026} & 2026 & \texttt{G18} & supplies a comparison \\
\texttt{\small CHEN\_\allowbreak{}2006} & \cite{CHEN_2006} & 2006 & \texttt{G18} & supplies a comparison \\
\texttt{\small LLL\_\allowbreak{}2006} & \cite{LLL_2006} & 2006 & \texttt{G18} & supplies a comparison \\
\texttt{\small MORNINGSTAR\_\allowbreak{}2025} & \cite{MORNINGSTAR_2025} & 2025 & \texttt{G18} & supplies a comparison \\
\texttt{\small MP\_\allowbreak{}1999} & \cite{MP_1999} & 1999 & \texttt{G18} & supplies a comparison \\
\texttt{\small AT\_\allowbreak{}2021\_\allowbreak{}SUN} & \cite{AT_2021_SUN} & 2021 & \texttt{G19} & supplies a comparison \\
\texttt{\small BB\_\allowbreak{}1983} & \cite{BB_1983} & 1983 & \texttt{G19} & supplies a comparison \\
\texttt{\small FLPS\_\allowbreak{}2026} & \cite{FLPS_2026} & 2026 & \texttt{G19} & supplies a comparison \\
\texttt{\small URZUA\_\allowbreak{}2019} & \cite{URZUA_2019} & 2019 & \texttt{G19} & supplies a comparison \\
\texttt{\small BORGA\_\allowbreak{}2024} & \cite{BORGA_2024} & 2024 & \texttt{G5} & supplies a comparison \\
\texttt{\small CAO\_\allowbreak{}2023} & \cite{CAO_2023} & 2023 & \texttt{G5} & supplies a comparison \\
\texttt{\small OBZ\_\allowbreak{}1985} & \cite{OBZ_1985} & 1985 & \texttt{G5} & supplies a comparison \\
\texttt{\small AT\_\allowbreak{}2021\_\allowbreak{}SUN} & \cite{AT_2021_SUN} & 2021 & \texttt{G6} & supplies a comparison \\
\texttt{\small CM\_\allowbreak{}2003} & \cite{CM_2003} & 2003 & \texttt{G6} & supplies a comparison \\
\texttt{\small BALAJI\_\allowbreak{}2026} & \cite{BALAJI_2026} & 2026 & \texttt{R2} & supplies a comparison \\
\texttt{\small CKS\_\allowbreak{}2021} & \cite{CKS_2021} & 2021 & \texttt{U1} & supplies a comparison \\
\texttt{\small DRS\_\allowbreak{}2021} & \cite{DRS_2021} & 2021 & \texttt{U1} & supplies a comparison \\
\texttt{\small KRS\_\allowbreak{}2023} & \cite{KRS_2023} & 2023 & \texttt{U1} & supplies a comparison \\
\texttt{\small SZH\_\allowbreak{}1997} & \cite{SZH_1997} & 1997 & \texttt{U1} & supplies a comparison \\
\texttt{\small CS\_\allowbreak{}2006} & \cite{CS_2006} & 2006 & \texttt{C7} & supplies a method \\
\texttt{\small JW\_\allowbreak{}2006} & \cite{JW_2006} & 2006 & \texttt{G17} & supplies a method \\
\texttt{\small UELTSCHI\_\allowbreak{}2004\_\allowbreak{}CLUSTER} & \cite{UELTSCHI_2004_CLUSTER} & 2004 & \texttt{G17} & supplies a method \\
\texttt{\small YAROTSKY\_\allowbreak{}2004\_\allowbreak{}GS} & \cite{YAROTSKY_2004_GS} & 2004 & \texttt{G17} & supplies a method \\
\texttt{\small BB\_\allowbreak{}1983} & \cite{BB_1983} & 1983 & \texttt{G18} & supplies a method \\
\texttt{\small JW\_\allowbreak{}2006} & \cite{JW_2006} & 2006 & \texttt{G18} & supplies a method \\
\texttt{\small NSY\_\allowbreak{}2019} & \cite{NSY_2019} & 2019 & \texttt{G18} & supplies a method \\
\texttt{\small YAROTSKY\_\allowbreak{}2004\_\allowbreak{}QP} & \cite{YAROTSKY_2004_QP} & 2004 & \texttt{G18} & supplies a method \\
\texttt{\small JW\_\allowbreak{}2006} & \cite{JW_2006} & 2006 & \texttt{G19} & supplies a method \\
\texttt{\small KLEIN\_\allowbreak{}ROSENBERGER\_\allowbreak{}2020\_\allowbreak{}WKB} & \cite{KLEIN_ROSENBERGER_2020_WKB} & 2020 & \texttt{G19} & supplies a method \\
\texttt{\small SCHIERHOLZ\_\allowbreak{}1988} & \cite{SCHIERHOLZ_1988} & 1988 & \texttt{G19} & supplies a method \\
\texttt{\small UELTSCHI\_\allowbreak{}2004\_\allowbreak{}CLUSTER} & \cite{UELTSCHI_2004_CLUSTER} & 2004 & \texttt{G19} & supplies a method \\
\texttt{\small YAROTSKY\_\allowbreak{}2004\_\allowbreak{}GS} & \cite{YAROTSKY_2004_GS} & 2004 & \texttt{G19} & supplies a method \\
\texttt{\small YAROTSKY\_\allowbreak{}2004\_\allowbreak{}QP} & \cite{YAROTSKY_2004_QP} & 2004 & \texttt{G19} & supplies a method \\
\texttt{\small JW\_\allowbreak{}2006} & \cite{JW_2006} & 2006 & \texttt{G20} & supplies a method \\
\texttt{\small JW\_\allowbreak{}2006} & \cite{JW_2006} & 2006 & \texttt{G22} & supplies a method \\
\texttt{\small JW\_\allowbreak{}2006} & \cite{JW_2006} & 2006 & \texttt{G23} & supplies a method \\
\texttt{\small MUNSTER\_\allowbreak{}1985\_\allowbreak{}TM} & \cite{MUNSTER_1985_TM} & 1985 & \texttt{G3} & supplies a method \\
\texttt{\small KRS\_\allowbreak{}2023} & \cite{KRS_2023} & 2023 & \texttt{R2} & supplies a method \\
\end{longtable}

\nocite{*}


